\section{What LLMs are good at}

The strengths that matter for GGM are those that fit the inputs described in Chapter 3: reports, spreadsheets in many formats, documentation and the judgment that ties them together.

\textbf{Reading unstructured material.} LLMs read text, tables and documents in whatever form they arrive and can extract the quantities a task needs, which is why most of the systems in Section 2.2 use them to turn reports and manuals into model input.
For economists, \citet{korinek2023generative} includes summarising and extracting from text among his use cases, and a later guide shows how agents can combine literature search, data retrieval and code for research \citep{korinek2025ai}.

\textbf{Applying written rules.} An instruction such as ``map each pipeline to the node of its entry point, and convert capacities to the model's units'' can be given in plain language, without first being written as a program.
This suits the rule-bound but irregular work of mapping infrastructure data, where the rules are clear but every source reports in its own way.

\textbf{Writing and repairing code.} Agentic coding is the best-measured agent capability.
SWE-bench tests whether a model can resolve real issues in open-source software repositories \citep{jimenez2024swe}, and it became the standard benchmark for agents that read code, change it and run tests.
Code written by a model can be read, tested and reused, which makes it a more checkable product than a number the model states directly.

\textbf{Revising under feedback.} Agents improve when they are told why an attempt failed.
Reflexion agents that write a short reflection on each failure and keep it for the next attempt outperform agents that simply retry \citep{shinn2023reflexion}, and the harnesses in Chapter 2 rely on the same mechanism when they return solver errors to the model.
Calibration, which is a sequence of runs, deviations and adjustments, has this structure.

\textbf{Explaining.} Because the model reasons in text, it can state why it proposes a change, and this rationale can be logged and reviewed.
The calibration agent in Section 2.3 attaches such a rationale to every iteration \citep{gomezguillen2026agentic}.

How good the best models are at expert work has been measured directly.
In GDPval, experts with on average fourteen years of experience in 44 occupations wrote realistic work tasks, and other experts compared the models' deliverables blind with those of human professionals; the authors conclude that the best models are approaching experts in deliverable quality \citep{patwardhan2025gdpval}.
The tasks were, however, specified in full and completed in one step, which is exactly what GGM's calibration is not.

Capability is also changing quickly.
METR measures the length of software tasks, in the time they take a skilled person, that a model completes with 50\% probability, and found that this time horizon has doubled about every seven months since 2019 \citep{kwa2025measuring}.
The measure has been widely used as evidence that agents can take on ever longer tasks, but its authors caution that a 50\% success rate is not a level at which work can safely be delegated, that the estimates are uncertain by about a factor of two, and that the tasks are cleaner than most real work \citep{metr2026clarifying}.
Any assessment of what agents can do in GGM's workflow is therefore a snapshot, and the division of work proposed in Chapter 5 should be revisited as models improve.
